% !TEX root = ../main.tex
\section{Explicación de cómo se almacenará la información en un archivo secuencial}

La información permanente vive en tres archivos binarios de registros de
tamaño fijo: \texttt{identidad.dat}, \texttt{academico.dat} y
\texttt{contacto.dat}, uno por cada struct definido en la sección de
diseño del registro. El tamaño fijo es lo que hace posible calcular la
posición física de cualquier registro sin recorrer el archivo:

\begin{lstlisting}[style=pseudocodigo]
offset = posicion_logica * sizeof(struct)
\end{lstlisting}

Con ese \texttt{offset}, \texttt{fseek} posiciona el cursor de lectura o
escritura directamente sobre el registro deseado, y \texttt{fread} o
\texttt{fwrite} lo leen o sobrescriben sin tocar el resto del archivo. En
ningún momento se carga el archivo completo a memoria: solo se traen a RAM
los registros puntuales que una operación necesita, siguiendo la posición
que entrega el mapa de posiciones correspondiente (ver sección de archivo
indexado).

Los tres archivos comparten la misma posición lógica por estudiante, de
modo que el registro en la posición $i$ de cada uno pertenece siempre al
mismo carné, sin necesidad de guardar una llave adicional dentro de cada
struct para relacionarlos entre sí.
